\lab{1}{Добыча корпуса документов}
\task
Подготовить корпус для следующих лабораторных: получить документы не менее
чем из двух источников, определить единицу документа, изучить разметку и
метаданные, выделить текст. Найти существующие средства поиска по источникам,
привести примеры запросов и ограничения выдачи. Определить количество
документов, объёмы исходных данных и текста, средние размеры документов.

\method
Корпус объединяет техническую документацию и архивные обсуждения Stack
Exchange. Для документации единицей является HTML-страница с собственным
URL и заголовком; для обсуждений - архивная ветка вопроса. Такое разделение
сохраняет контекст задачи и ответов, но приводит к неодинаковой длине
документов. Источники включают Linux Kernel, man7, cppreference, проекты
LLVM, GNU, FreeBSD и восемь сайтов Stack Exchange.

Документация поступает из экспортированной базы робота. Для большого
числа обсуждений используется набор \code{common-pile/stackexchange}
с Hugging Face \cite{dataset}. Зафиксирована ревизия
\code{5ec3aa2e65b7c61b7c18f1ff18ccbc4d3279a23c};
происхождение и авторство не подменяются результатом собственного обхода.
Изначально порог длины обсуждения составлял 1000 слов; начиная с архива
\code{00024_se.jsonl.gz} для Stack Overflow он снижен до 850 слов.
Страницы документации этим порогом не ограничены.

HTML обрабатывается профилем конкретного сайта: удаляются оформление,
скрипты и навигация, сохраняется основной текст, включая примеры кода.
Для повторного извлечения в исходном экспорте остаются HTML-файлы.
Большой корпус хранится компактно, без отдельного текстового файла на
каждый документ. Формат \code{corpus.bin} начинается сигнатурой
\code{IRDOC001}, далее следуют записи:
\begin{center}
\begin{tabular}{lll}\toprule
Поле & Размер & Содержание\\\midrule
\code{url_length} & 4 байта & Длина URL в байтах\\
\code{title_length} & 4 байта & Длина заголовка в байтах\\
\code{text_length} & 8 байт & Длина текста в байтах\\
URL, заголовок, текст & переменный & UTF-8, подряд без разделителей\\\bottomrule
\end{tabular}
\end{center}
Целые записываются little-endian. Манифест JSONL хранит идентификаторы,
источник, URL, заголовок и смещение текста; сводная статистика хранится
отдельно. Это позволяет читать корпус последовательно и проверять
целостность независимо от индексатора.

\journal
Сначала реализован экспорт из MongoDB и извлечение текста по профилям
сайтов. Затем добавлены импорт архивных обсуждений, проверка архивов и
контрольные точки: прерванный импорт продолжается с последнего сохранённого
состояния. Большое количество мелких файлов заменено упакованным потоком.
Для защиты от неполного экспорта анализатор сверяет смещения и объёмы
с манифестом. Главная проблема неоднородного корпуса - разные понятия
«сырого документа»: для HTML это исходная страница, для архива -
распакованная JSON-строка. В статистике эта договорённость записана явно.

\testing
Проверяются сигнатура потока, длины полей, отсутствие NUL внутри текста,
корректность UTF-8, монотонность смещений и выход за границы файла.
Число записей и сумма длин текстов должны совпасть с итоговыми агрегатами.
Отдельно следует проверять извлечение на страницах разных источников:
осмысленный заголовок, наличие основного текста, сохранение примеров кода,
отсутствие меню. Проверка количества документов сама по себе не выявляет
дублирующиеся тексты, поэтому дедупликация по содержимому остаётся отдельной
задачей. Статистика включает число документов и объёмы
по каждому источнику, что позволяет независимо сверять итоговые суммы.

\results
\begin{table}[H]\centering
\begin{tabular}{lr}\toprule
Показатель & Значение\\\midrule
Количество документов & 1\,679\,971\\
Число источников в статистике & 22\\
Исходные данные, байт & 20\,327\,391\,393\\
Извлечённый текст, байт & 17\,236\,388\,223\\
Упакованный корпус, байт & 17\,433\,701\,772\\
Слов в статистике экспорта & 2\,237\,727\,857\\
Средний исходный размер, байт & 12\,099,85\\
Средний объём текста, байт & 10\,259,93\\
Среднее число слов & 1332,00\\\bottomrule
\end{tabular}
\caption{Статистика основного корпуса. «Слова» экспортёра и токены индексатора считаются разными правилами.}
\end{table}
\begin{longtable}{lrr}
\caption{Распределение документов и текста по источникам.}\\
\toprule Источник & Документов & Текст, MiB\\\midrule\endfirsthead
\toprule Источник & Документов & Текст, MiB\\\midrule\endhead
clang & 144 & 4,56\\
freebsd-docs & 183 & 12,60\\
linux-kernel & 3\,516 & 37,41\\
cppreference & 6\,838 & 20,41\\
gcc & 314 & 2,17\\
llvm-libc & 133 & 0,45\\
llvm-libcxx & 38 & 0,44\\
lld & 14 & 0,06\\
lldb & 9\,598 & 85,21\\
llvm & 314 & 6,49\\
man7 & 5\,473 & 22,59\\
binutils & 4 & 0,02\\
gdb & 1 & 0,02\\
glibc & 306 & 1,20\\
codereview.stackexchange.com & 33\,980 & 461,43\\
softwareengineering.stackexchange.com & 21\,456 & 240,22\\
unix.stackexchange.com & 18\,921 & 184,43\\
dba.stackexchange.com & 9\,995 & 101,76\\
serverfault.com & 19\,463 & 192,03\\
askubuntu.com & 22\,181 & 242,55\\
superuser.com & 26\,517 & 242,14\\
stackoverflow.com & 1\,500\,582 & 14579,69\\
\bottomrule
\end{longtable}


Stack Overflow составляет около 89,3\% документов. Поэтому даже большое
число источников не означает равномерное покрытие тем. Суммарная статистика
также не означает, что каждый документ содержит несколько тысяч слов:
среди справочных страниц есть короткие описания API.

Для внешнего поиска подходят запросы
\code{site:docs.kernel.org memory allocation},
\code{site:man7.org pthread_create},
\code{site:en.cppreference.com std::vector} и поиск
\code{[c++] memory allocation} по Stack Overflow.
Они проверяют соответственно тему, точное имя функции и технический
идентификатор. Ограничения такого сравнения: внешняя система индексирует
другую версию сайтов, может включать отсутствующие в архиве страницы и
ранжирует их по собственным сигналам. Сохранённого снимка внешней выдачи
нет, поэтому этим примерам не приписываются измеренные Precision или NDCG.

\conclusions
Получен тематический корпус, который можно обрабатывать потоком и который
превышает миллион документов. Объём извлечённого текста составляет около
16,05 GiB. Основные ограничения - преобладание Stack Overflow,
неодинаковая длина документов и возможные содержательные дубликаты.
Улучшения: дедупликация по нормализованному содержимому, стратифицированная
оценка по источникам и отдельный контроль качества извлечения кода.
